\chapter*{General Conclusion and Perspectives}


% This conclusion is unnumbered, if you want it numbered, you can remove the * from above and remove the line below, so it becomes a chapter, then add sections.

\addcontentsline{toc}{chapter}{General Conclusion and Perspectives} %adds to the table of contents
\markboth{General Conclusion and Perspectives}{General Conclusion and Perspectives} % running header of the next pages

\label{chap:General Conclusion}

This internship, carried out at the Cognitus Global Delivery and Support Center in Tangier from June 15 to August 15, 2026, was my first immersion in the SAP ecosystem. Its goal was to bring together two worlds that recruitment usually keeps apart: a modern self-service experience for external candidates and the governance of HR data in the enterprise system. The answer to this challenge is Talently, a recruitment and employee administration platform built as a side-by-side extension on SAP BTP.

The work started with an upskilling phase on SAP Learning (BTP100, S4D400, and S4D437), followed by the analysis of the recruitment process, the specification of the requirements, and the design of the solution. Talently relies on three components. The backend, developed with ABAP RAP on the SAP BTP ABAP Environment, holds all business data and rules in six business objects, from the organizational structure to the employee master record, and exposes them through OData V4 services; the HR cockpit is generated by SAP Fiori Elements from the annotations of the CDS views. The candidate portal, built with React and Clerk, lets candidates find offers, apply, follow their progress, and complete their onboarding form. Finally, a FastAPI microservice analyzes each resume with GPT-4o-mini, enriches it with the candidate's GitHub projects, and scores it against the rubric of the role: the language model judges, but the final score is computed in Python.

The results show that the objectives set in Chapter~\ref{chap:project_context_objectives} have been reached. The end-to-end scenario covered the whole process, from the creation of the organizational units to the administration of the new employee, and every group of functional requirements was validated. On the SAP BTP trial system, all SAP operations answer in less than half a second, and the hiring of an accepted candidate creates the employee master record and its child records in a single transaction, in about 0.4~s and without any manual re-entry. The 19 tests of the microservice and the four ABAP Unit tests pass, and the security checks confirmed that write requests with an outdated ETag, without an \texttt{If-Match} header, or without a CSRF token are rejected. The AI engine ranked three fictitious profiles of known level in the expected order in every run, for about 0.003~USD per resume. Its evaluation also revealed a fairness defect: personal details were still sent to the scoring model. Once the scoring input was anonymized, the difference between a resume and its copy with another name, gender, city, and university was no longer significant. This experience showed the value of testing an AI component as rigorously as the rest of the code.

The prototype also has limitations. It runs on an SAP BTP trial account: the portal reaches the OData services through a development proxy that reuses an SAP session, and the AI microservice runs on the development machine. The AI score is stable but not fully deterministic, since the evaluation of the open source contributions varies from one run to another, and the engine was only tested on a small set of fictitious resumes and with a single role rubric. Finally, the candidates must share their resume through a public Google Drive link, which is practical for a prototype but not suited to personal data in production.

These limitations open several perspectives:

\begin{itemize}
    \item \textbf{Production deployment:} replace the development proxy with a server-side integration that calls the Web API \texttt{ZAPI\_HCM\_RECRUIT} through its communication arrangement, and host the portal and the Docker image of the AI microservice in the cloud.
    \item \textbf{Integration with the SAP HR core:} replicate the employees created by Talently to the standard employee master data of SAP S/4HANA or SAP SuccessFactors Employee Central through their released APIs.
    \item \textbf{Reliability and fairness of the AI engine:} use the median of three scoring calls, always add the GitHub profile found by the pipeline to the parsed resume, define rubrics for other roles, and compare the rankings of the engine with those of recruiters on a larger set of anonymized resumes.
    \item \textbf{Data protection:} store the resumes in a secure document service instead of Google Drive, and define the consent, retention, and deletion rules required by the applicable data protection laws, such as the Moroccan law 09-08 and the GDPR.
    \item \textbf{Recruiter experience:} send automatic notifications to the candidates at each pipeline step, and add recruitment indicators such as the time-to-hire and the conversion rate of each step through analytical CDS queries.
\end{itemize}

On a personal level, this internship was a rich learning experience. I discovered the ERP world and the SAP technologies, from ABAP Cloud to RAP, CDS, OData, and Fiori Elements, and I learned to apply SAP's Clean Core principle to a concrete business case. Building a complete product, from its charter and planning to its tests, taught me to organize my work in sprints, to document my decisions, and to measure results instead of assuming them. I also learned that integrating artificial intelligence into an enterprise process requires as much care for fairness and transparency as for performance. These skills will be a solid basis for the rest of my engineering studies and for my future career.
